% Déclaration sur l'utilisation de l'IA / Declaration on the use of AI
%
% Requested by the jury (J3). The text is the letter to the jury in ai-usage-wzz.md,
% with the letter framing removed and the language cleaned up.
\ifthenelse{\equal{\Langue}{english}}{
	\chapter*{DECLARATION ON THE USE OF AI}\thispagestyle{headings}
	\addcontentsline{toc}{compteur}{DECLARATION ON THE USE OF AI}
}{
	\chapter*{DÉCLARATION SUR L'UTILISATION DE L'IA}\thispagestyle{headings}
	\addcontentsline{toc}{compteur}{DÉCLARATION SUR L'UTILISATION DE L'IA}
}
This page describes how I used AI tools in the work presented in this thesis on improving the
canonicalization time of \texttt{CVXPY}.

As mentioned in the acknowledgements, I used AI tools, in particular Claude Code from
Anthropic, in almost all parts of the research process. The usage varied greatly from one part
to another, so each is described in its own section below.

\section*{Coding}

By coding, I mean the open-source code contributions that have been, or will be, added to the
\texttt{CVXPY} and \texttt{SparseDiffEngine} repositories. I consider this code more important
than the experiment code described below, since it can affect many users.

For coding, I first started with a plan of what I wanted to add and iterated with Claude Code
until I was satisfied. Claude Code then implemented the change, wrote new tests and verified
that everything worked. After that, I opened a pull request and did a second, more thorough code
review, which sometimes included stylistic cleanups. After a few more iterations on the code,
one or several of the other maintainers reviewed the pull request and eventually merged it.
Although Claude Code was used at almost every stage, from making a plan to writing the code and
sometimes even acting as a reviewer, the process still depended heavily on my own ideas and
vision for the software.

In my experience, this is more or less standard practice in software engineering today. I have
thought a lot about AI usage in open source, and it may be unfortunate that it has gone this
way, but it is at least the practice that the two main packages I contribute to have recently
adopted.

\section*{Experiments}

In contrast to coding, by experiments I mean more ``throw-away'' code written to test ideas or
improvements I had in mind. Experiments sometimes became more serious, and then fell into the
category of coding.

For experiments, I was much less strict about code quality than in the section above. What I
verified thoroughly is that the experiments ran the right versions of the software and measured
the right thing. Beyond that, I left a lot of leeway to Claude Code, which could produce messy
and duplicated code. I am currently reorganizing the experiments to keep only the final ones
that appear in this thesis, and will make them reproducible, easier to read, and of a quality
similar to the code described in the section above.

\section*{Research}

By research, I mean specifically understanding related work for the literature review and
solidifying the story of the contribution. The main ideas of this thesis are my own, but many
of the improvements between an initial idea and a completed thesis were supported by AI tools.
To understand other works, Claude Code can coordinate agents that search for papers and
documentation online, process them, and report back to me. I then receive a condensed summary,
focused on what I am looking for, which I verify myself. I have not blindly trusted AI tools,
because they can make mistakes, but I have used them as a search aggregator, and they have been
tremendously helpful in showing me where to look.

To polish the story of the contribution, I used AI tools as a conversational companion that
could check the validity of some ideas and make them more concrete. I also benefited greatly
from discussions with Daniel Cederberg, the other developer of \texttt{SparseDiffEngine}, and
with Oscar Dowson, the main developer of \texttt{JuMP}, which is covered in the literature
review. My advisor, Prof.\ Antoine Lesage-Landry, gave me feedback on asking the right questions
and discussing the right topics.

It is hard to quantify the contribution of AI tools in this regard, but I am confident that
most of the main ideas and the shape of the contribution came from us, with AI tools helping
only occasionally to verify, solidify and question our approach.

\section*{Writing}

Finally, I used Claude Code to assist with the prose of this thesis. A large part of this usage
was improving the English and the phrasing of the points I wanted to make. The workflow was as
follows. I first wrote down on paper the high-level ideas I wanted in a section or paragraph.
I then tried to form coherent sentences in a long prompt to Claude Code, which edited and
organized them before adding them to the \texttt{.tex} file.

Claude Code also gave immediate feedback on my writing, and sometimes suggested ways to
elaborate on how I expressed my thoughts. I always proofread and iterated many times with
Claude Code before committing to a particular phrasing.

A smaller part of my AI usage in writing was generating starting sentences that gave me an idea
of what the final text could look like. This was especially helpful when I was stuck without any
idea of what to write. Although the generated sentences were not always used, they gave me some
editing momentum. In these cases, I acted more as a reviewer, using my taste and my vision of
what the paragraph should say to direct the writing of Claude Code.

Overall, I take full responsibility for every claim in this thesis. For the literature review,
I verified the \texttt{JuMP} section with Oscar Dowson, who is the expert on the subject, and
searched through much of the history and documentation of \texttt{AMPL} to validate my claims
about how it works. On \texttt{CVXPY}, I am confident about the exposition, since I have
substantial experience with its internal codebase and have verified what I describe with other
maintainers, including Steven Diamond, the original developer of \texttt{CVXPY}. Although I used
AI extensively in my writing, the final product is the result of many choices of my own, from
the initial organization of the thesis to proofreading, many times, every sentence that the AI
edited or generated.

\section*{Miscellaneous}

Other uses worth mentioning are the figures, and the formatting and editing of the thesis
document, all in \LaTeX. I delegated these tasks almost completely to AI tools, which were
surprisingly good at them. I verified that the resulting PDF matches what I expected or
wanted, but looked much less at the \texttt{.tex} source itself, especially for the figures.
